% Français 9115, batch 18 — Gallica views 341–360.
% Pass: Opus 5.5 (claude-opus-5-5), 2026-09-22 — first pass, unchecked against the views by a human.
%
% What the twenty views hold. One piece, begun before the batch and going
% on after it: a fair copy, in French, of Sophie Germain's critique of
% Euler's elastic lamina touched by a stylet (« stilet ») at an interior
% point — the same argument, under the same paragraph numbers, as the
% worked draft transcribed in batch 14 (views 262–276), here without its
% hesitations. The hand is the regular, upright fair hand of the Latin
% Euler copy (views 144–206), with the same « u » for alpha and « v » for
% gamma; the corrections are few (three struck groups, three or four words
% added above the line or in the margin). Each leaf is written on one side only, carries a
% parenthesised page number at the head, struck through, and a number in
% ink at the top right, and is mounted on a guard, overlapping the leaf
% before it.
%
% Ten views carry writing: 342, 344, 346, 348, 350, 352, 354, 356, 358,
% and 359, which is not a new page but the leaf of view 358 photographed a
% second time (same text, same blots, same « 175 »), and is not
% transcribed. The nine odd views 341, 343 … 357 and view 360 are the blank
% backs of the leaves, showing the writing of the leaf through the paper,
% reversed, with a strip of the previous leaf's back at the top where the
% mounting overlaps; they were looked at in full and are not transcribed.
% No microfilm target, no printed matter.
%
%   v. 342  (12) opens mid-calculation: equations I and II of the
%           preceding page added to the two members, β and β' replaced by
%           their values, and the final equation for a lamina with both
%           ends free and the stylet at the fraction λ of its length; § 9,
%           the case λ = 1/2 and its two solutions — the first makes the
%           two halves vibrate as two laminae of length a/2, each free at
%           one end and appuyée at the other.
%   v. 344  (13) that condition checked against Euler's « cas I § 5 »;
%           the second solution (halves fixed at the stylet) rejected
%           « comme j'ai montré qu'il falloit le faire pour celle qui lui
%           est analogue dans le cas examiné par Euler »; § 10, where the
%           analogy with the supported ends breaks down: there any aliquot
%           fraction λ gives two laminae of lengths λa and (1−λ)a.
%   v. 346  (14) for free ends, only λ = 1/2 will do; § 11, the proof that
%           the two conditions cannot vanish together, reduced to
%           2 sin λω {(e^{(1−2λ)ω} + e^{(2λ−1)ω}) sin ω − (e^ω − e^{−ω}) cos ω}.
%   v. 348  (15) sin λω = 0 excluded; the second factor compared with
%           Euler's end condition gives 2 = ±(e^{(1−2λ)ω} + e^{(2λ−1)ω}),
%           impossible unless ω = 0 or λ = 1/2; § 12, the same results for
%           both ends fixed, with its final equation.
%   v. 350  (16) the fixed-ends case at λ = 1/2; why only the middle
%           allows the lamina to be taken as two separate laminae — equal
%           parts between a node and the stylet on either side.
%   v. 352  (17) what remains to be proved; § 14 (no § 13 is written),
%           Euler's ordinate at a point L at λa and at λ ∓ λ' (« V. M.
%           d'Euler § 21 »), their sum set to zero, reduced to
%           {e^{λω} ∓ e^{(1−λ)ω}}{e^{λ'ω} + e^{−λ'ω} − 2 cos λ'ω} = 0; the
%           fixed-ends variant begun.
%   v. 354  (18) the fixed-ends variant finished: only λ = 1/2, an odd
%           number of nodes, one of them at the middle. Then, in the first
%           person, « Je ne crois pas que l'on ait encore fait cette
%           remarque de l'inégalité des parties comprises entre deux nœuds
%           de vibrations »: Chladni (« V. traité d'acoustique P. 94 § 71 »)
%           and Giordano Ricati, « Delle vibrazioni de' cilindri », first
%           volume of the Memorie of the Società italiana, « (Voyez. § XXX
%           P 503) » with the paragraph left as three crosses, as in the
%           draft.
%   v. 356  (19) what Ricati's values show (the part at the end less than
%           half the next, the parts growing towards the middle) and that
%           he drew no conclusion; § 15, the supported ends, where the
%           parts are equal because sin(λ−λ')ω + sin(λ+λ')ω =
%           2 sin λω cos λ'ω vanishes with sin λω.
%   v. 358  (20) § 16, the analogy between the final equations: the four
%           final equations (both ends appuyées § 5, free § 8, fixed § 12,
%           one fixed and one appuyée), and the list of Euler's six end
%           conditions with his section numbers 5, 23, 30, 35, 39, 43. The
%           page ends there; view 361 (outside this batch) opens § 17,
%           page (21), ink number 176.
%
% Seams. View 342 continues view 340 (page (11), ink 166, « Ainsi
% l'équation III se réduit… », ending on equation IV), which belongs to
% batch 17: the equations I, II, III and the values of β and β' that 342
% calls « cidessus » are there. View 358 ends at the end of a list; view
% 361 opens § 17 on a new page, so the piece runs on into batch 19 without
% a sentence cut.
%
% The numbers on the leaves, and why no \folio{} is written. Two series,
% one per written leaf, none on the blank backs. (1) At the head, centre,
% in parentheses and struck through with a horizontal stroke (a cross at
% v. 352): (12), (13), (14), (15), (16), (17), (18), (19), (20) at views
% 342 to 358 — the fair copy's own pagination, continuing (11) at v. 340;
% the figure at v. 346 is overwritten and read (14) from the run. (2) At
% the top right, in ink, in the larger rounder hand of the running series
% of this volume: 167, 168, 169, 170, 171, 172, 173, 174, 175 at views 342
% to 358, one per leaf without a gap, continuing 166 at v. 340 and going on
% to 176 at v. 361; 175 is repeated on the second exposure, v. 359. As in
% every batch so far, the series occupies the place of the Bibliothèque's
% foliation but is penned, not pencilled, so no \folio{} is written. (3)
% Her paragraph numbers, inside the first line: 9 (v. 342), 10 (v. 344),
% 11 (v. 346), 12 (v. 348), 14 (v. 352), 15 (v. 356), 16 (v. 358); no 13
% appears. Every number above was read on its view; none is inferred.
% None of these leaves can be Laubenbacher and Pengelley's « Manuscript C »
% (ff. 348r–349r): the subject is the elastic lamina, and the ink series
% here runs 166–176.
%
% Notation. Hers, kept. λ is the fraction of the length a at which the
% stylet stands (a small hooked letter, read λ as on the Euler copy), λ'
% the fraction to the neighbouring nodes, its accent written like an acute
% on the letter; ω the angle of Euler's cases; α, β, γ, δ and α', β', γ',
% δ' Euler's coefficients (α written « u », γ written « v »). « sin »,
% « cos », « tang » are set \text{sin}, \text{cos}, \text{tang}; the stop
% she writes after them (« sin.λω ») is irregular and is not recorded stop
% by stop. Braces and square brackets are hers; where she leaves a brace or
% a parenthesis unclosed (views 342, 358) it is left unclosed. Where the
% leaf and the calculation disagree the leaf stands, with a note: v. 344
% (« et sin(1−λ)ω » without « = 0 »), v. 348 (an exponent (2λ−1) without
% ω), v. 352 (∓ where the equation it comes from has ±), v. 354
% (e^{λω} = ± e^{−λω} where (1−λ)ω is expected). Spelling is hers and is
% not flagged: « coëfficiens », « extremités », « dela », « demême »,
% « aulieu », « àprésent », « parconséquent », « desorte », « cidessus »,
% « précédens », « différens », « fesant », « appuiée » beside « appuyée »,
% « égals », « etres », « rejetter », « appercevoir », « tems », « Ricati »,
% and agreements such as « deux lames séparés », « extremités appuyés ».
%
% What could not be read, and what is offered with doubt. No \ill{} in the
% batch. \uncertain{} stands three times: « Voye. » (v. 344), for a
% contraction of « Voyez »; « 21 » in « § 21 » (v. 352), the second figure
% merging with the closing parenthesis; « 94 » in Chladni's page number
% (v. 354), which may be « g4 », i.e. « pag 4 » as the draft was read.
%
% Nothing here was checked against any published edition of Euler's memoir,
% of Chladni, of Riccati or of Sophie Germain's papers.
\documentclass[11pt,a4paper]{article}
% The preamble shared by every transcription of the mathematicians' manuscripts in Gallica.
%
% It rests on one idea: **uncertainty compounds, it does not resolve.** A
% transcription that smooths over a doubtful word has destroyed the only thing
% it was made for — a reader must be able to tell, at every point, what was on
% the page and what was supplied. The apparatus macros below are therefore the
% critical apparatus, not decoration.
%
% The same commands are recognised by `scripts/render.mjs`, which produces the
% site's reading view, and by `scripts/tei.mjs`. A macro added here must be
% added there too, or rendering fails loudly — which is the wanted behaviour.
%
% Comments here are English, like the rest of the repository. The documents
% this preamble sets are French, and that is deliberate: see the skills.

\ProvidesPackage{germain}[2026/09/15 Transcriptions of mathematicians' manuscripts in Gallica]

\usepackage{amsmath,amssymb,amsthm}
\usepackage{mathrsfs}
% Kept from the parent preamble so the renderer's diagram support stays
% available; these manuscripts draw no commutative diagrams, and nothing here uses it.
\usepackage{tikz-cd}
\usepackage[english,french]{babel}
\usepackage{fontspec}

% --- Greek ------------------------------------------------------------------
% The text face stays fontspec's default, Latin Modern, which has no Greek:
% XeTeX drops every glyph it lacks with a line in the log and nothing on the
% page, and the Greek words the manuscripts quote vanished from the PDFs. So
% the Greek and Greek Extended blocks get a XeTeX character class of their own,
% and entering or leaving a run of it switches the family to CMU Serif — the
% same Computer Modern design, with polytonic Greek — and back. Only the
% family changes: a Greek word in \textit or \textbf keeps its shape and series.
% The transitions to and from every other class carry the tokens that class
% has towards an ordinary letter, so French spacing before « ; » or inside
% « » still holds after a Greek word. No grouping, so no brace or \endgroup
% can come out unbalanced; maths is left alone, since XeTeX inserts nothing
% there. Kept identical in `exercices.sty`.
\makeatletter
\newfontfamily\germainGreekfont{cmun}[
  Extension=.otf, NFSSFamily=germaingreek,
  UprightFont=*rm, ItalicFont=*ti, SlantedFont=*sl,
  BoldFont=*bx, BoldItalicFont=*bi, BoldSlantedFont=*bl]
\newXeTeXintercharclass\germain@greek
\def\germain@greekrange#1#2{%
  \@tempcnta=#1\relax
  \loop
    \XeTeXcharclass\@tempcnta=\germain@greek
    \ifnum\@tempcnta<#2\advance\@tempcnta\@ne\repeat}
\germain@greekrange{"0370}{"03FF}
\germain@greekrange{"1F00}{"1FFF}
\def\germain@greekprev{\rmdefault}
\def\germain@greekin{%
  \edef\germain@greekprev{\f@family}\fontfamily{germaingreek}\selectfont}
\def\germain@greekout{\fontfamily{\germain@greekprev}\selectfont}
\def\germain@greekjoin{%
  \XeTeXinterchartokenstate=\@ne
  \@tempcnta=\z@
  \loop
    \ifnum\@tempcnta=\germain@greek\else
      \XeTeXinterchartoks\germain@greek\@tempcnta=\expandafter{%
        \expandafter\germain@greekout\the\XeTeXinterchartoks\z@\@tempcnta}%
      \XeTeXinterchartoks\@tempcnta\germain@greek=\expandafter{%
        \the\XeTeXinterchartoks\@tempcnta\z@\germain@greekin}%
    \fi
    \ifnum\@tempcnta<4095\advance\@tempcnta\@ne\repeat}
% babel-french sets its own classes, and saves and restores the interchar
% state, each time a language is selected: join again after it does.
\addto\extrasfrench{\germain@greekjoin}
\addto\extrasenglish{\germain@greekjoin}
\AtBeginDocument{\germain@greekjoin}
\makeatother

\usepackage{geometry}
\geometry{a4paper,margin=25mm,marginparwidth=28mm}
\usepackage{marginnote}
\usepackage{xcolor}
\usepackage[normalem]{ulem}
\setlength{\emergencystretch}{3em}
\usepackage{hyperref}
\hypersetup{colorlinks=true,linkcolor=black,urlcolor=black,pdfborder={0 0 0}}

% --- Batch metadata ---------------------------------------------------------
% Taken as they stand by the rendering script, which shows them at the head of
% the reading view. They fix what the file claims to cover. The macro names are
% the parent project's, so that the scripts stay shared; \folder here means the
% volume's slug — `fr-9115` — and \pages counts Gallica views.

\newcommand{\folder}[1]{\gdef\grothFolder{#1}}
\newcommand{\batch}[1]{\gdef\grothBatch{#1}}
\newcommand{\pages}[2]{\gdef\grothFirst{#1}\gdef\grothLast{#2}}
\newcommand{\foldertitle}[1]{\gdef\grothFoldertitle{#1}}

% The holder's own shelfmark, printed as they write it: « Français 9115 ».
\newcommand{\shelfmark}[1]{\gdef\grothShelfmark{#1}}

% The dating, copied verbatim from the holder's catalogue — never inferred
% here. « XIXe siècle » is the BnF's; a pass that dates a leaf from its content
% says so in a \note{}, as its own inference.
\newcommand{\dating}[1]{\gdef\grothDating{#1}}

% The status notice, printed in the title block and repeated as a diagonal
% watermark on every page of the PDF. The manuscripts are in the public
% domain; what this line declares is not a right but a quality — an
% unverified first pass by a machine. The skills write the line; a file
% without it gets no watermark, so leaving it out is visible in the output.
\newcommand{\watermark}[1]{\gdef\grothWatermark{#1}}

\gdef\grothFolder{?}\gdef\grothBatch{}\gdef\grothFirst{?}\gdef\grothLast{?}%
\gdef\grothFoldertitle{}\gdef\grothShelfmark{}\gdef\grothDating{}\gdef\grothWatermark{}

% --- The résumé -------------------------------------------------------------
\newenvironment{resume}
  {\small\setlength{\parskip}{0.45em}}
  {\par\medskip\hrule\medskip}

% --- Keywords ---------------------------------------------------------------
% In English, closing the modernised reading's résumé: the modern vocabulary
% under which to search for what the leaves construct. The manifest extracts
% them to tag the volume — this line is the single source of the tags.
\newcommand{\keywords}[1]{\par\smallskip\noindent
  {\footnotesize\textsc{Keywords} --- \textit{#1}}\par}

% --- The critical apparatus -------------------------------------------------

% The Gallica view number, in the margin. This is the anchor that lets a
% disagreement over a formula be settled by looking at *one* image rather than
% a volume of 750. The site uses it to turn the facsimile as the transcription
% is scrolled. A view is one scanned image: usually one side of a leaf.
\newcommand{\page}[1]{%
  \marginnote{\footnotesize\color{gray!70}\textsf{v.\,#1}}%
  \label{page:#1}\ignorespaces}

% The BnF's pencilled foliation, where the leaf carries it and a pass can read
% it: « 198r ». This is what the literature cites — « ff. 198r–208v » — and
% the one line that turns a citation into a view. Recorded, never inferred:
% a folio a pass cannot read is not written.
\newcommand{\folio}[1]{%
  \marginnote{\footnotesize\color{gray!50}\textsf{f.\,#1}}\ignorespaces}

% The modernised reading's coarser counterpart to \page{}: which views a
% section is drawn from.
\newcommand{\pagerange}[2]{%
  \marginnote{\footnotesize\color{gray!70}\textsf{v.\,#1--#2}}%
  \ignorespaces}

% Illegible: never guessed.
\newcommand{\ill}{\textcolor{red!60!black}{[\,\ldots\,]}}

% A doubtful reading: the word is given, and flagged as uncertain.
\newcommand{\uncertain}[1]{\uline{#1}}

% An editorial addition — a missing letter, an implied word.
\newcommand{\add}[1]{\textcolor{blue!50!black}{[#1]}}

% Struck out by the writer. What was crossed out is part of the document.
\newcommand{\struck}[1]{\sout{#1}}

% The transcriber's note, on the state of the page or the mathematics.
\newcommand{\note}[1]{\par\smallskip{\small\color{gray!60!black}\textit{[#1]}}\par\smallskip}

% A marginal note of hers, distinct from the body of the text.
\newcommand{\marginal}[1]{\par\smallskip\noindent\hspace{1em}%
  {\small\color{gray!60!black}#1}\par\smallskip}

% --- Title ------------------------------------------------------------------
% The title block is French, like the documents themselves.

\AddToHook{shipout/background}{%
  \ifx\grothWatermark\empty\else
    \begin{tikzpicture}[remember picture, overlay]
      \node[rotate=52, scale=3.4, align=center, text=gray!18,
            font=\sffamily\bfseries]
        at (current page.center) {\grothWatermark};
    \end{tikzpicture}%
  \fi}

\AtBeginDocument{%
  \begin{center}
    {\large\bfseries Manuscrits de mathématiciens (Gallica) ---
      \ifx\grothShelfmark\empty\grothFolder\else\grothShelfmark\fi}\\[2pt]
    {\itshape\grothFoldertitle}\\[4pt]
    {\small vues \grothFirst--\grothLast
      \ifx\grothBatch\empty\else\ (lot \grothBatch)\fi}\\[2pt]
    \ifx\grothDating\empty\else
      {\small\color{gray!60!black}Datation du catalogue : \grothDating}\\[2pt]
    \fi
    \ifx\grothWatermark\empty\else
      {\small\bfseries\color{red!45!gray}\grothWatermark}\\[2pt]
    \fi
    {\footnotesize\color{gray!60!black}Transcription établie d'après les images IIIF de Gallica.
     Source gallica.bnf.fr / Bibliothèque nationale de France.\\
     Les lectures incertaines sont soulignées, les illisibles notées [\,\ldots\,],
     les ajouts entre crochets ; \textsf{v.} numérote les vues de Gallica,
     \textsf{f.} la foliotation au crayon de la BnF.}
  \end{center}
  \vspace{1em}}

\endinput


\folder{fr-9115}
\shelfmark{Français 9115}
\batch{18}
\pages{341}{360}
\dating{1801-1900}
\watermark{Lecture automatique\\non vérifiée}
\foldertitle{Papiers de Sophie GERMAIN. — Recueil de dissertations et problèmes mathématiques et physiques.}

\begin{document}

\note{Numérotation des feuillets. Chaque feuillet écrit porte deux nombres,
et les dos blancs n'en portent aucun : au milieu de la tête, entre
parenthèses et barré, le numéro de page de cette mise au net, (12) à (20)
aux vues 342 à 358 ; en haut à droite, à l'encre, d'une écriture plus
ronde, 167 à 175 aux mêmes vues, un par feuillet et sans lacune. Cette
seconde série occupe la place de la foliotation de la Bibliothèque, mais
elle est tracée à la plume et non au crayon, et elle prolonge la série à
l'encre relevée dans les lots précédents de ce volume ; aucune
foliotation au crayon n'a été lue ici, et aucun « f. » n'est donc porté.
Tous ces nombres ont été lus sur la vue ; aucun n'est déduit. Les vues
341, 343, 345, 347, 349, 351, 353, 355, 357 et 360 sont les dos blancs
des feuillets ; la vue 359 reprend le feuillet de la vue 358 ; aucune
n'est transcrite.}

\page{342}

en ajoutant I.~$\alpha e^{\lambda\omega} + \beta e^{-\lambda\omega} + \gamma\,\text{sin}\,\lambda\omega + \delta\,\text{cos}\,\lambda\omega = 0$ au premier membre,

II.~$\alpha' e^{\lambda\omega} + \beta' e^{-\lambda\omega} + \gamma'\,\text{sin}\,\lambda\omega + \delta'\,\text{cos}\,\lambda\omega = 0$ au second, et divisant par 2; il en resulte
$\alpha e^{\lambda\omega} + \beta e^{-\lambda\omega} = \alpha' e^{\lambda\omega} + \beta' e^{-\lambda\omega}$. cette dernière équation devient, en y mettant pour
$\beta$ et $\beta'$ les valeurs cidessus
\[ \frac{\alpha e^{\lambda\omega}\{e^{-\lambda\omega} - \text{sin}\,\lambda\omega + \text{cos}\,\lambda\omega\} - \alpha e^{-\lambda\omega}\{e^{\lambda\omega} + \text{sin}\,\lambda\omega + \text{cos}\,\lambda\omega\}}{e^{-\lambda\omega} - \text{sin}\,\lambda\omega + \text{cos}\,\lambda\omega} \]
\[ = \frac{\begin{array}{c}\alpha' e^{(\lambda-1)\omega}\{e^{(1-\lambda)\omega} + \text{cos}\,(1-\lambda)\omega + \text{sin}\,(1-\lambda)\omega\} \\ {} - \alpha' e^{(1-\lambda)\omega}\{e^{(\lambda-1)\omega} + \text{cos}\,(1-\lambda)\omega - \text{sin}\,(1-\lambda)\omega\}\end{array}}{e^{-\omega}\{e^{(1-\lambda)\omega} + \text{cos}\,(1-\lambda)\omega + \text{sin}\,(1-\lambda)\omega\}} \]
ou,
\begin{align*}
&\frac{\alpha\{(e^{\lambda\omega} - e^{-\lambda\omega})\,\text{cos}\,\lambda\omega - (e^{\lambda\omega} + e^{-\lambda\omega})\,\text{sin}\,\lambda\omega\}}{e^{-\lambda\omega} - \text{sin}\,\lambda\omega + \text{cos}\,\lambda\omega} \\
&= \frac{\alpha'\{(e^{(1-\lambda)\omega} + e^{(\lambda-1)\omega})\,\text{sin}\,(1-\lambda)\omega - (e^{(1-\lambda)\omega} - e^{(\lambda-1)\omega})\,\text{cos}\,(1-\lambda)\omega\}}{e^{-\omega}\{e^{(1-\lambda)\omega} + \text{cos}\,(1-\lambda)\omega + \text{sin}\,(1-\lambda)\omega}
\end{align*}
D'ou on tire,
\begin{align*}
&\frac{\alpha e^{-\omega}\{e^{(1-\lambda)\omega} + \text{cos}\,(1-\lambda)\omega + \text{sin}\,(1-\lambda)\omega\}}{\alpha'\{e^{-\lambda\omega} - \text{sin}\,\lambda\omega + \text{cos}\,\lambda\omega\}} \\
&= \frac{(e^{(1-\lambda)\omega} + e^{(\lambda-1)\omega})\,\text{sin}\,(1-\lambda)\omega - (e^{(1-\lambda)\omega} - e^{(\lambda-1)\omega})\,\text{cos}\,(1-\lambda)\omega}{(e^{\lambda\omega} - e^{-\lambda\omega})\,\text{cos}\,\lambda\omega - (e^{\lambda\omega} + e^{-\lambda\omega})\,\text{sin}\,\lambda\omega}
\end{align*}

Il ne s'agit plus àprésent que demettre à la place du premier membre de cette équation
sa valeur déduite de l'équation III, pour conclure
\[ \{2 + (e^{(1-\lambda)\omega} + e^{(\lambda-1)\omega}\,\text{cos}\,(1-\lambda)\omega\}\{(e^{\lambda\omega} + e^{-\lambda\omega})\,\text{sin}\,\lambda\omega - (e^{\lambda\omega} - e^{-\lambda\omega})\,\text{cos}\,\lambda\omega\} \]
\[ + \{2 + (e^{\lambda\omega} + e^{-\lambda\omega})\,\text{cos}\,\lambda\omega\}\{(e^{(1-\lambda)\omega} + e^{(\lambda-1)\omega})\,\text{sin}\,(1-\lambda)\omega - (e^{(1-\lambda)\omega} - e^{(\lambda-1)\omega})\,\text{cos}\,(1-\lambda)\omega = 0 \]
cette dernière équation est l'équation finale, dont il faut tirer la valeur
de l'angle $\omega$.

9. Il est aisé devoir que, dans le cas ou $\lambda = \frac{1}{2}$, l'équation finale se réduit à
\[ 2\{2 + (e^{\frac{1}{2}\omega} + e^{-\frac{1}{2}\omega})\,\text{cos}\,\tfrac{1}{2}\omega\}\{(e^{\frac{1}{2}\omega} + e^{-\frac{1}{2}\omega})\,\text{sin}\,\tfrac{1}{2}\omega - (e^{\frac{1}{2}\omega} - e^{-\frac{1}{2}\omega})\,\text{cos}\,\tfrac{1}{2}\omega\} = 0 \]
et que cette équation est susceptible de deux solutions, savoir:
\[ (e^{\frac{1}{2}\omega} + e^{-\frac{1}{2}\omega})\,\text{sin}\,\tfrac{1}{2}\omega - (e^{\frac{1}{2}\omega} - e^{-\frac{1}{2}\omega})\,\text{cos}\,\tfrac{1}{2}\omega = 0 \quad\text{et}\quad 2 + (e^{\frac{1}{2}\omega} + e^{-\frac{1}{2}\omega})\,\text{cos}\,\tfrac{1}{2}\omega = 0. \]
La première est relative au cas où le stilet est consideré comme simplement
appuyé aumilieu de la lame dont les extremités sont libres; et elle montre que
les deux moitiés de cette lame doivent alors vibrer comme deux lames séparés
de longueur $\frac{1}{2}a$ qui auroient chacune une de leurs extremités libre, et l'autre
appuyé; car la condition $\frac{e^{\frac{1}{2}\omega} - e^{-\frac{1}{2}\omega}}{e^{\frac{1}{2}\omega} + e^{-\frac{1}{2}\omega}} = \text{tang}\,\tfrac{1}{2}\omega$ est celle qui convient au
mouvement d'une telle lame. Cette solution s'accorde au reste avec l'équation
dépendante de l'état des extremités dela lame primitive,
\note{En tête, au milieu, « (12) » entre parenthèses, barré d'un trait
horizontal ; en haut à droite « 167 », à l'encre, d'une écriture plus ronde.
Le texte fait suite à la vue 340 (« (11) », « 166 »), qui n'est pas de ce
lot : les équations I, II, III et les valeurs de $\beta$ et $\beta'$
« cidessus » y sont. La lettre écrite ici $\lambda$ est un petit signe
crochu, proche d'un « A » ou d'un « r », que le passage sur la copie
d'Euler a déjà rendu par $\lambda$ ; $\alpha$ est tracé comme un « u » et
$\gamma$ comme un « v ». Dans l'égalité qui suit « ou, », l'accolade du
dénominateur de droite n'est pas refermée ; dans les deux lignes de
l'équation finale, une parenthèse puis une accolade fermantes manquent sur
le feuillet ; on les laisse manquer. Au premier membre de la première
égalité, le second terme porte $\alpha$ et non $\alpha'$, comme on lit.
« séparés » et « appuyé » sont ainsi accordés sur le feuillet. Le
paragraphe porte le numéro 9, écrit dans la ligne.}

\page{344}

car on a (\uncertain{Voye.} M. d'Euler, cas I § 5) $\text{cos}\,\omega = \frac{2e^{\omega}}{e^{2\omega} + 1}$, d'où l'on tire, en prenant $\text{sin}\,\omega$
positif comme il convient ici, $2\,\text{cos}^2\,\tfrac{1}{2}\omega = \frac{(e^{\omega} + 1)^2}{e^{2\omega} + 1}$, $2\,\text{sin}\,\tfrac{1}{2}\omega\,\text{cos}\,\tfrac{1}{2}\omega = \frac{e^{2\omega} - 1}{e^{2\omega} + 1}$, $\text{tang}\,\tfrac{1}{2}\omega = \frac{e^{\omega} - 1}{e^{\omega} + 1}$,
qui est précisément la condition résultante de la première solution.

À l'égard de la seconde solution, elle donne $\frac{2}{e^{\frac{1}{2}\omega} + e^{-\frac{1}{2}\omega}} = -\text{cos}\,\tfrac{1}{2}\omega$: ainsi
elle meneroit à conclure que la lame primitive peut encore être considerée
comme separée en deux autres lames de longueur $\frac{1}{2}a$ qui auroient chacune
une de leurs extremités fixée au point d'application du stilet et l'autre libre.
mais la condition $\frac{2}{e^{\frac{1}{2}\omega} + e^{-\frac{1}{2}\omega}} = -\text{cos}\,\tfrac{1}{2}\omega$ ne peut s'accorder avec l'équation
$\text{cos}\,\omega = \frac{2e^{\omega}}{e^{2\omega} + 1}$ dépendante de l'état des extremités, que dans la seule
supposition $\omega = 0$. Il faut donc rejetter cette seconde solution, comme j'ai montré
qu'il falloit le faire pour celle qui lui est analogue dans le cas examiné par
Euler

10 Jusqu'ici il existe la plus grande analogie entre le cas où les deux
extremités sont appuyés et celui ou les deux extremités sont libres; car, de
part et d'autre, la supposition $\lambda = \frac{1}{2}$, donne des solutions semblables,
quant a la considération de l'effet de l'application du stilet: mais
malheureusement un examen plus approfondi du cas présent, fait appercevoir
de grandes différences. En effet, dans le cas des extremités appuyés, la
première solution relative à la supposition $\lambda = \frac{1}{2}$ n'est qu'un cas particulier
dela solution générale, dans laquelle $\lambda$ peut être une fraction quelconque
aliquote de l'unité; desorte que, quelque soit celle que l'on adopte, la lame
primitive peut toujours être considerée comme separée en deux autres lames
des longueurs $\lambda a$ et $(1-\lambda)a$, qui auroient, l'une et l'autre, leurs deux extremités
appuyés; et cette conclusion résulte de ceque les équations $\text{sin}\,\lambda\omega = 0$ et
$\text{sin}\,(1-\lambda)\omega$, ont lieu en meme tems pour de telles valeurs de $\lambda$, et de ceque
ces valeurs sont les seules pour lesquelles le mouvement dela lame soit
possible,
\note{En tête, au milieu, « (13) » barré d'un trait ; en haut à droite
« 168 », à l'encre. « car » est écrit un peu au-dessus de la ligne, en
tête ; « que l'on » est ajouté au-dessus de la ligne, entre « celle » et
« adopte ». La référence lit « Voye. », peut-être pour « Voyez » ; le
nom est tracé avec l'E qui ressemble à un L, comme ailleurs dans ce
volume. Le « m » de « examen » porte un trait suscrit. « et
$\text{sin}\,(1-\lambda)\omega$ » est écrit sans « $= 0$ ». « appuyés »,
« séparée », « considerée » sont accordés ainsi sur le feuillet. Le
paragraphe porte le numéro 10, dans la ligne. La page s'achève sur
« possible, » au quart inférieur du feuillet.}

\page{346}

Au contraire, dans le cas présent des deux extremités libres, les équations
\[ \frac{e^{(1-\lambda)\omega} + e^{(\lambda-1)\omega}}{e^{(1-\lambda)\omega} - e^{(\lambda-1)\omega}}\,\text{sin}\,(1-\lambda)\omega - \text{cos}\,(1-\lambda)\omega = 0 \quad\text{et}\quad \frac{e^{\lambda\omega} + e^{-\lambda\omega}}{e^{\lambda\omega} - e^{-\lambda\omega}}\,\text{sin}\,\lambda\omega - \text{cos}\,\lambda\omega = 0 \]
ne peuvent avoir lieu en même tems, lorsqu'on a égard à la condition qui résulte de l'état
des extremités dela lame primitive, que dans le seul cas où $\lambda = \frac{1}{2}$. ces équations
peuvent etre mises sous la forme $\text{tang}\,(1-\lambda)\omega = \frac{e^{(1-\lambda)\omega} - e^{(\lambda-1)\omega}}{e^{(1-\lambda)\omega} + e^{(\lambda-1)\omega}}$, $\text{tang}\,\lambda\omega = \frac{e^{\lambda\omega} - e^{-\lambda\omega}}{e^{\lambda\omega} + e^{-\lambda\omega}}$
qui montre qu'elles ne sont autre chose que les conditions du mouvement de
deux lames des longueurs $\lambda a$ et $(1-\lambda)a$, dont une des extremités seroit appuyée, et
l'autre libre. Ainsi on est forcé de conclure que le cas où le stilet est appuyé
au milieu dela lame dont les deux extremités sont libres, est le seul où cette
lame puisse etre consideré comme separé en deux autres lames, qui auroient
chacune une de leurs extremités appuyée au point d'application du stilet, et
l'autre libre.

11 Maintenant je vais prouver qu'en effet les deux quantités $\frac{e^{(1-\lambda)\omega} + e^{(\lambda-1)\omega}}{e^{(1-\lambda)\omega} - e^{(\lambda-1)\omega}}\,\text{sin}\,(1-\lambda)\omega - \text{cos}\,(1-\lambda)\omega$
et $\frac{e^{\lambda\omega} + e^{-\lambda\omega}}{e^{\lambda\omega} - e^{-\lambda\omega}}\,\text{sin}\,\lambda\omega - \text{cos}\,\lambda\omega$, ne peuvent être supposées en même tems égales à zéro.
pour cela j'emploie les transformations suivantes, dans lesquelles les valeurs tirées de l'équation
$\frac{e^{\lambda\omega} + e^{-\lambda\omega}}{e^{\lambda\omega} - e^{-\lambda\omega}}\,\text{sin}\,\lambda\omega - \text{cos}\,\lambda\omega = 0$ sont introduites dans celle ci $\frac{e^{(1-\lambda)\omega} + e^{(\lambda-1)\omega}}{e^{(1-\lambda)\omega} - e^{(\lambda-1)\omega}}\,\text{sin}\,(1-\lambda)\omega - \text{cos}\,(1-\lambda)\omega = 0$
\[ \frac{e^{(1-\lambda)\omega} + e^{(\lambda-1)\omega}}{e^{(1-\lambda)\omega} - e^{(\lambda-1)\omega}}\,\text{sin}\,(1-\lambda)\omega - \text{cos}\,(1-\lambda)\omega \]
\[ = \frac{e^{(1-\lambda)\omega} + e^{(\lambda-1)\omega}}{e^{(1-\lambda)\omega} - e^{(\lambda-1)\omega}}\{\text{sin}\,\omega\,\text{cos}\,\lambda\omega - \text{sin}\,\lambda\omega\,\text{cos}\,\omega\} - \{\text{cos}\,\omega\,\text{cos}\,\lambda\omega + \text{sin}\,\omega\,\text{sin}\,\lambda\omega\} \]
\begin{align*}
&= \frac{e^{(1-\lambda)\omega} + e^{(\lambda-1)\omega}}{e^{(1-\lambda)\omega} - e^{(\lambda-1)\omega}}\Bigl\{\frac{e^{\lambda\omega} + e^{-\lambda\omega}}{e^{\lambda\omega} - e^{-\lambda\omega}}\cdot\text{sin}\,\lambda\omega\,\text{sin}\,\omega - \text{sin}\,\lambda\omega\,\text{cos}\,\omega\Bigr\} \\
&- \Bigl\{\frac{e^{\lambda\omega} + e^{-\lambda\omega}}{e^{\lambda\omega} - e^{-\lambda\omega}}\,\text{sin}\,\lambda\omega\,\text{cos}\,\omega + \text{sin}\,\omega\,\text{sin}\,\lambda\omega\Bigr\} = 0
\end{align*}
et encore, en multipliant par $(e^{\lambda\omega} - e^{-\lambda\omega})(e^{(1-\lambda)\omega} - e^{(\lambda-1)\omega})$,
\begin{align*}
0 = \text{sin}\,\lambda\omega\,\bigl\{&(e^{\omega} + e^{(2\lambda-1)\omega} + e^{(1-2\lambda)\omega} + e^{-\omega})\,\text{sin}\,\omega - (e^{\omega} + e^{(2\lambda-1)\omega} - e^{(1-2\lambda)\omega} - e^{-\omega})\,\text{cos}\,\omega \\
&- (e^{\omega} - e^{(2\lambda-1)\omega} + e^{(1-2\lambda)\omega} - e^{-\omega})\,\text{cos}\,\omega - (e^{\omega} - e^{(2\lambda-1)\omega} - e^{(1-2\lambda)\omega} + e^{-\omega})\,\text{sin}\,\omega\bigr\}
\end{align*}
\[ = 2\,\text{sin}\,\lambda\omega\,\{(e^{(1-2\lambda)\omega} + e^{(2\lambda-1)\omega})\,\text{sin}\,\omega - (e^{\omega} - e^{-\omega})\,\text{cos}\,\omega\} \]
Elles montrent que, si les équations dont il s'agit pouvoient etres satisfaites, il faudroit
\note{En tête, au milieu, un nombre entre parenthèses, barré et
surchargé, qui paraît être « (14) » ; en haut à droite « 169 », à
l'encre. Le texte suit la vue 344 sans lacune. Le paragraphe 11 porte son
numéro dans la marge, dans la ligne. « peuvent », dans « ne peuvent être
supposées », est récrit sur un mot plus court ; un petit signe à trois
points se lit au-dessus de « j'emploie », sans renvoi visible. Dans la
première égalité, le signe « $-$ » du premier dénominateur est repassé et
empâté. Le dernier facteur de la grande accolade est écrit sur deux
lignes, la seconde alignée à droite sous la première. « etres » pour
« être », « consideré comme separé » : ainsi sur le feuillet. La page
s'arrête sur « il faudroit » ; la suite est à la vue 348.}

\page{348}

que l'un ou l'autre des facteurs $\text{sin}\,\lambda\omega$ et $(e^{(1-2\lambda)\omega} + e^{(2\lambda-1)})\,\text{sin}\,\omega - (e^{\omega} - e^{-\omega})\,\text{cos}\,\omega$ fut égal
à zéro; mais, à moins de faire $\lambda$ ou $\omega = 0$, $\text{sin}\,\lambda\omega$ ne peut etre supposé ici égal à zéro,
parceque l'équation $\text{sin}\,\lambda\omega = 0$ appartient au mouvement d'une lame de longueur $\lambda a$
dont les deux extremités sont appuyés, et qu'il s'agit présentement d'une lame de
même longueur, dont une seule extrémité est censé appuyée, et l'autre libre.
ainsi il ne reste à examiner que l'équation $(e^{(1-2\lambda)\omega} + e^{(2\lambda-1)\omega})\,\text{sin}\,\omega - (e^{\omega} - e^{-\omega})\,\text{cos}\,\omega = 0$
ou $\text{tang}\,\omega = \frac{e^{\omega} - e^{-\omega}}{e^{(1-2\lambda)\omega} + e^{(2\lambda-1)\omega}}$; et la comparaison dela valeur de $\text{tang}\,\omega$, tirée
des équations $\text{cos}\,\omega = \frac{2}{e^{\omega} + e^{-\omega}}$ $\text{sin}\,\omega = \pm\frac{e^{\omega} - e^{-\omega}}{e^{\omega} + e^{-\omega}}$ relatives à la considération de l'état
des extremités dela lame primitive, donne $\text{tang}\,\omega = \pm\frac{e^{\omega} - e^{-\omega}}{2} = \frac{e^{\omega} - e^{-\omega}}{e^{(1-2\lambda)\omega} + e^{(2\lambda-1)\omega}}$
ou $2 = \pm(e^{(1-2\lambda)\omega} + e^{(2\lambda-1)\omega})$. il est clair que le cas relatif au signe $-$ doit être
rejetté, parceque 2 et $e^{(1-2\lambda)\omega} + e^{(2\lambda-1)\omega}$ sont essentiellement des quantités positives; et
on voit encore avec facilité que $e^{(1-2\lambda)\omega} + e^{(2\lambda-1)\omega}$ est toujours $>$ que 2,
excepté pour le cas $\omega = 0$ qui appartient au repos, et pour lavaleur de $\lambda$, $\lambda = \frac{1}{2}$.

12. On obtient des résultats semblables pour le cas où les extrémités dela lame,
aulieu d'être libres, sont l'une et l'autre fixées; car, si on suppose qu'une
telle lame soit soumise à l'action d'un stilet dans un point quelconque de
sa longueur, on parvient au moyen de calculs que je supprime ici, à cause
de leur parfaite analogie avec les précédens, à l'équation finale
\[ \Bigl\{\frac{2}{e^{\lambda\omega} + e^{-\lambda\omega}} - \text{cos}\,\lambda\omega\Bigr\}\Bigl\{\frac{e^{(1-\lambda)\omega} + e^{(\lambda-1)\omega}}{e^{(1-\lambda)\omega} - e^{(\lambda-1)\omega}}\cdot\text{sin}\,(1-\lambda)\omega - \text{cos}\,(1-\lambda)\omega\Bigr\} \]
\[ + \Bigl\{\frac{2}{e^{(1-\lambda)\omega} + e^{(\lambda-1)\omega}} - \text{cos}\,(1-\lambda)\omega\Bigr\}\Bigl\{\frac{e^{\lambda\omega} + e^{-\lambda\omega}}{e^{\lambda\omega} - e^{-\lambda\omega}}\,\text{sin}\,\lambda\omega - \text{cos}\,\lambda\omega\Bigr\} = 0 \]
delaquelle on tire les valeurs de l'angle $\omega$. Et on trouve, comme dans
le cas précédent qu'en fesant $\lambda = \frac{1}{2}$, on a pour solution de cette équation
la condition $\frac{e^{\frac{1}{2}\omega} + e^{-\frac{1}{2}\omega}}{e^{\frac{1}{2}\omega} - e^{-\frac{1}{2}\omega}}\cdot\text{sin}\,\tfrac{1}{2}\omega - \text{cos}\,\tfrac{1}{2}\omega = 0$, qui mène à conclure que
\note{En tête, au milieu, « (15) » barré d'un trait ; en haut à droite
« 170 », à l'encre. Le texte suit la vue 346 sans lacune. Au premier
facteur de la première ligne, le second exposant est écrit $(2\lambda-1)$,
sans $\omega$ ; on le laisse ainsi. « supposé », « appuyés », « censé
appuyée » : ainsi sur le feuillet. Deux petits traits obliques se lisent
au-dessus de « le cas », avant « précédent », sans renvoi visible. Le
paragraphe 12 porte son numéro dans la ligne. La page s'arrête sur « qui
mène à conclure que » ; la suite est à la vue 350.}

\page{350}

la lame primitive peut se mouvoir, comme si elle étoit séparée en deux
autres lames demoitié moins longues, qui auroient, chacune, une de leurs
extremités appuyée au point d'application du stilet, et l'autre fixée.
On trouve encore, par le calcul déja employé pour le cas précédent que
les équations $\text{tang}\,(1-\lambda)\omega = \frac{e^{(1-\lambda)\omega} - e^{(\lambda-1)\omega}}{e^{(1-\lambda)\omega} + e^{(\lambda-1)\omega}}$, $\text{tang}\,\lambda\omega = \frac{e^{\lambda\omega} - e^{-\lambda\omega}}{e^{\lambda\omega} + e^{-\lambda\omega}}$ ne peuvent
se concilier avec la condition dépendante de l'état des extremités pour
aucune autre valeur de $\lambda$ que $\lambda = \frac{1}{2}$.

Il semble d'abord étonnant que, dans le cas où $\lambda = \frac{1}{2}$, la lame dont les deux
extremités sont fixés ou libres, puisse etre consideré comme reellement separée
en deux autres lames, qui auroient une de leurs extremités appuyée au point
d'application du stilet tandis que la même supposition n'est plus applicable à aucun autre
cas

Voici comment on peut se rendre compte de cette singularité: si plusieurs lames
dont les deux extremités sont libres ou fixées, rendent le même son, et sont pourtant
de longueurs différentes, le nombre des nœuds de vibration relatif à chacune d'elles
sera différent; mais l'étendue des parties vibrantes comprises entre les extremités et le nœud
le plus voisin de chacune d'elles, restera lemême; et il en sera demême, si les lames
supposées ont une seule de leurs extremités libre ou fixée, et l'autre appuiée.
Cela posé, si on considère la lame élastique vibrante dont les deux extremités
sont libres ou fixées, comme separée en deux parties réellement indépendantes l'une
de l'autre, et l'une et l'autre appuyés au point d'application du stilet, \struck{il faudra}
il faudra que les deux parties vibrantes comprises entre ce point et les nœuds de
vibrations les plus voisins prises audelà et endeçà de ce point, soient egales
entr'elles; car elles appartiendront dans cette hypothèse à deux lames de
longueurs différentes, dont une des extremités sera libre ou fixée et l'autre appuyée.
\note{En tête, au milieu, « (16) » barré d'un trait ; en haut à droite
« 171 », à l'encre. Le texte suit la vue 348 sans lacune. « du stilet »
est ajouté au-dessus de la ligne, après « d'application » ; « il
faudra », écrit au bout d'une ligne, y est barré et récrit en tête de la
suivante. Une tache d'encre précède « le même son ». Le mot rendu
« nœuds » est écrit avec une ligature qui se lit presque « naud ».
« fixés », « consideré », « appuyés », « appuiée » : ainsi sur le
feuillet. La page s'achève sur « l'autre appuyée. », au bas du
feuillet.}

\page{352}

et qui, par la nature dela question, doivent rendre le même son, il ne reste
donc plus, pour éclaircir la difficulté, qu'à prouver qu'excepté le cas où il y a
un nœud au milieu dela longueur dela lame, il n'y a aucun point de repos tel
que les ordonnées prises à distances egales de part et d'autre de ce point, soient égales
entr'elles, et que parconséquent les parties vibrantes comprises entre le point d'applica-tion
du stilet consideré comme simple point d'appui et les nœuds le plus voisins
pris audelà et endeçà de ce point, ne peuvent être égales entr'elles que dans le
seul cas où le stilet est supposé appliqué au milieu dela lame dont il s'agit.

14. Pour démontrer cette proposition en s'occupant d'abord du cas où les deux
extremités dela lame, \struck{il faut se rappeller} sont libres, il faut se rappeller que,
pour un point quelconque L situé aux distances $a\lambda$ et $(1-\lambda)a$ des extremités dela lame,
la supposition $y = 0$ qui résulte de ceque le stilet est appuyé dans ce point,
donne l'équation $e^{\lambda\omega} \pm e^{(1-\lambda)\omega} + (1 \mp e^{\omega})\,\text{sin}\,\lambda\omega + (1 \pm e^{\omega})\,\text{cos}\,\lambda\omega = 0$ (V. M. d'Euler § \uncertain{21}).
on voit aisément que les ordonnées des points situés àdroite et àladistance $\lambda' a$ du
point L, seront
\[ \ldots\ldots\{e^{(\lambda-\lambda')\omega} \pm e^{(1-\lambda+\lambda')\omega} + (1 \mp e^{\omega})\,\text{sin}\,(\lambda-\lambda')\omega + (1 \pm e^{\omega})\,\text{cos}\,(\lambda-\lambda')\omega\} = y \]
\[ \ldots\ldots\{e^{(\lambda+\lambda')\omega} \pm e^{(1-\lambda-\lambda')\omega} + (1 \mp e^{\omega})\,\text{sin}\,(\lambda+\lambda')\omega + (1 \pm e^{\omega})\,\text{cos}\,(\lambda+\lambda')\omega\} = -y \]
ces ordonnées sont de signes contraires, parceque, si l'une est située au dessus
de l'axe, l'autre doit l'être au dessous du même axe. si elles sont égales entr'elles,
leur somme sera donc nulle, et on aura
\[ 0 = \{e^{\lambda\omega} \pm e^{(1-\lambda)\omega}\}\{e^{\lambda'\omega} + e^{-\lambda'\omega}\} + (1 \mp e^{\omega}).\,2\,\text{sin}\,\lambda\omega\,\text{cos}\,\lambda'\omega + (1 \pm e^{\omega})\,2\,\text{cos}\,\lambda\omega\,\text{cos}\,\lambda'\omega \]
et, à cause de $e^{\lambda\omega} \mp e^{(1-\lambda)\omega} = -\{(1 \mp e^{\omega})\,\text{sin}\,\lambda\omega + (1 \pm e^{\omega})\,\text{cos}\,\lambda\omega\}$,
cette condition se réduira à $\{e^{\lambda\omega} \mp e^{(1-\lambda)\omega}\}\{e^{\lambda'\omega} + e^{-\lambda'\omega} - 2\,\text{cos}\,\lambda'\omega\} = 0$.

À l'égard du cas où les deux extremités sont supposées fixées, la seule
différence entre les valeurs des coëfficiens $\alpha$, $\beta$, $\gamma$ et $\delta$ comparées à celles
qui sont relatives à l'hypothèse des extremités libres, est que les deux
\note{En tête, au milieu, « (17) », barré de deux traits croisés ; en
haut à droite « 172 », à l'encre. Le texte suit la vue 350 sans lacune.
Le paragraphe est numéroté 14, dans la ligne : aucun numéro 13 ne se lit
entre le 12 de la vue 348 et celui-ci. « égales », au bout de la
quatrième ligne, est surchargé et empâté. « il faut se rappeller » est
barré dans la ligne et récrit après « sont libres ». Dans la référence à
Euler, le chiffre qui suit le 2 se confond avec la parenthèse fermante :
« § 21 » ou « § 2 ». L'accent de $\lambda'$ est tracé comme un accent
aigu sur la lettre. Dans les deux dernières égalités, le signe entre
$e^{\lambda\omega}$ et $e^{(1-\lambda)\omega}$ est nettement $\mp$, alors
que l'équation de départ porte $\pm$ ; on garde les signes du feuillet.
Les quatre coefficients sont écrits « u, $\beta$, v et $\delta$ ». La
page s'arrête sur « est que les deux » ; la suite est à la vue 354.}

\page{354}

dernieres sont prises avec des signes contraires; et il en résulte seulement que l'on a
\[ e^{\lambda\omega} \pm e^{(1-\lambda)\omega} - \{(1 \mp e^{\omega})\,\text{sin}\,\lambda\omega + (1 \pm e^{\omega})\,\text{cos}\,\lambda\omega\} = 0 \]
\[ \ldots\ldots\{e^{(\lambda-\lambda')\omega} \pm e^{(1-\lambda+\lambda')\omega} - [(1 \mp e^{\omega})\,\text{sin}\,(\lambda-\lambda')\omega + (1 \pm e^{\omega})\,\text{cos}\,(\lambda-\lambda')\omega]\} = y \]
\[ \ldots\ldots\{e^{(\lambda+\lambda')\omega} \pm e^{(1-\lambda-\lambda')\omega} - [(1 \mp e^{\omega})\,\text{sin}\,(\lambda+\lambda')\omega + (1 \pm e^{\omega})\,\text{cos}\,(\lambda+\lambda')\omega]\} = -y \]
aulieu des valeurs précédentes; cequi n'empêche pas de conclure
\[ \{e^{\lambda\omega} \mp e^{(1-\lambda)\omega}\}\{e^{\lambda'\omega} + e^{-\lambda'\omega} - 2\,\text{cos}\,\lambda'\omega\} = 0 \]
et, pour satisfaire à cette condition, il faut que l'un ou l'autre des facteurs
$e^{\lambda\omega} \mp e^{(1-\lambda)\omega}$. et $e^{\lambda'\omega} + e^{-\lambda'\omega} - 2\,\text{cos}\,\lambda'\omega$ puisse être égal à zéro: mais on voit aisément
que l'équation $e^{\lambda'\omega} + e^{-\lambda'\omega} = 2\,\text{cos}\,\lambda'\omega$ ne peut être admise pour aucune valeur
de $\lambda'$ différente de zéro; ainsi il ne reste à examiner que l'équation $e^{\lambda\omega} = \pm e^{-\lambda\omega}$;
et il est évident que la seule valeur de $\lambda$ qui convienne, est $\lambda = \frac{1}{2}$ qui répond
au cas où il faut prendre le signe supérieur c'est adire à celui où il y a un
nombre impair de nœuds de vibrations, et parconséquent où il en existe un placé
au milieu dela lame.

Je ne crois pas que l'on ait encore fait cette remarque de l'inégalité des parties
comprises entre deux nœuds de vibrations dans les lames élastiques dont les deux extremités
sont libres ou fixées; car M\textsuperscript{r} Chladni se contente de dire, en parlant dela lame
dont les deux extremités sont libres, qu'une partie située à une extremité est
àpeuprès moitié moins longue qu'une partie comprise entre deux nœuds (V. traité d'acoustique
P. \uncertain{94} § 71); cequi semble supposer que toutes les parties soumises à cette dernière
condition sont entr'elles d'égale longueur. Cependant M\textsuperscript{r} Giordano Ricati, dans
son memoire \emph{Delle vibrazioni de' cilindri} inséré dans le premier
volume des mémoires dela société italienne, donne pour l'angle qui a été
désigné ici par $\lambda\omega$, des valeurs qui peuvent servir à confirmer la remarque
dont il s'agit pour le cas des extremités libres, seul cas qu'il ait traité
(Voyez. § XXX P 503)
\note{En tête, au milieu, « (18) » barré d'un trait ; en haut à droite
« 173 », à l'encre. Le texte suit la vue 352 sans lacune. Dans « que
l'équation $e^{\lambda\omega} = \pm e^{-\lambda\omega}$ », l'exposant de
droite est court et serré et se lit $-\lambda\omega$ ; le raisonnement
attendrait $(1-\lambda)\omega$ : on garde ce qu'on lit. Dans la
référence à Chladni, le nombre de la page se lit « g4 » ou « 94 », la
première figure ayant une queue descendante ; le brouillon de ce passage,
à la vue 272, a été lu « pag 4 §§ 71 ». Dans le renvoi final, trois
petites croix tiennent la place du numéro de paragraphe, comme au
brouillon de la vue 272. Le titre italien est écrit en caractères plus
gros que le reste de la ligne ; il est rendu en italique. Le nom est
écrit « Ricati ». Le feuillet s'achève sur ce renvoi, au bas de la
page.}

\page{356}

En effet, on voit, en comparant ces valeurs, qu'une partie située à l'extrémité, est
moins grande que la moitié dela partie comprise entre les deux nœuds les plus
voisins de cette extrémité; que la partie comprise entre le second nœud et le
troisième, est plus grande que celle qui est située entre le premier et le second;
que pareillement encore la partie comprise entre \struck{deux nœuds} le troisième et le
quatrième nœud, est plus grande que toutes les précédentes; desorte qu'en général
une partie comprise entre deux nœuds, est d'autant plus étendue qu'elle est
plus éloignée de l'extrémité et plus voisine du milieu dela lame, quoique
pourtant les différences soient beaucoup plus grandes entre les parties qui
sont près des extrémités, qu'entre celles qui s'en éloignent d'avantage.
Mais M\textsuperscript{r} Ricati s'est contenté de donner les valeurs dont il est question,
sans les comparer entr'elles, et sans conclure parconséquent l'inégalité des
parties vibrantes comprises entre deux nœuds.

15. Au contraire, lorsqu'il s'agit d'une lame élastique vibrante dont
les deux extremités sont simplement appuyés, les parties vibrantes comprises
entre deux nœuds sont égals entr'elles; car si, après avoir supposé que
le point d'application du stilet équivalant à un point d'appui soit
placé à la distance $\lambda a$ d'une des extremités, on veut qu'il existe
deux autres nœuds aux distances $(\lambda-\lambda')a$ et $(\lambda+\lambda')a$ de la même extrémité
et parconséquent situés l'un et l'autre à la distance $\lambda' a$ du point d'application
du stilet et de part et d'autre de ce même point; ces différentes conditions
seront exprimées par les équations $\text{sin}\,\lambda\omega = 0$ $\text{sin}\,(\lambda-\lambda')\omega = 0$, $\text{sin}\,(\lambda+\lambda')\omega = 0$;
et, en ajoutant les deux dernières ensemble, on trouvera la condition $2\,\text{sin}\,\lambda\omega\,\text{cos}\,\lambda'\omega = 0$
condition qui est évidemment remplie en vertu dela première et indépendamment
de la valeur de $\lambda'$. Ainsi on voit àprésent par quelle raison le cas des
deux extremités appuyés est le seul pour lequel on puisse avoir une solution
générale de l'équation finale trouvée par la méthode d'Euler.
\note{En tête, au milieu, « (19) » barré ; en haut à droite « 174 », à
l'encre. Le texte suit la vue 354 sans lacune. « moins » est écrit plus
petit, en retrait à gauche, au début de la deuxième ligne ; « pareillement »
est d'une écriture plus fine, dans la ligne, peut-être sur un mot gratté ;
« deux nœuds » est barré après « entre ». Plusieurs mots sont surchargés
sans que la lecture en soit changée (« l'extrémité » à la première ligne,
« des » devant « parties vibrantes »). Des taches d'encre marquent la
page près de « point » et de « première ». Le paragraphe 15 porte son
numéro dans la ligne. « égals », « appuyés » : ainsi sur le feuillet. La
page s'achève sur « la méthode d'Euler. », au bas du feuillet.}

\page{358}

16. On peut remarquer qu'il existe la plus grande analogie entre les diverses
équations finales relatives aux différens cas où on suppose le stilet appliqué
dans un point quelconque de la longueur des lames élastiques vibrantes dont
les extremités sont appuyées libres ou fixées, et aussi entre les facteurs de
ces équations et les conditions par lesquelles l'angle $\omega$ est déterminé
lorsqu'on a simplement égard à l'état des extremités des mêmes lames.
afin que l'on puisse saisir plus facilement ces rapports, je vais écrire ici
ces différentes équations, en y joignant encore l'équation finale que j'ai trouvée
pour le cas où le stilet est censé appliqué à la lame dont une des extremités
est fixée et l'autre appuyée

équation finale pour le cas des deux extremités appuyées: § 5
\begin{align*}
&\text{sin}\,\lambda\omega\,\Bigl\{\frac{e^{(1-\lambda)\omega} + e^{(\lambda-1)\omega}}{e^{(1-\lambda)\omega} - e^{(\lambda-1)\omega}}\,\text{sin}\,(1-\lambda)\omega - \text{cos}\,(1-\lambda)\omega \\
&+ \text{sin}\,(1-\lambda)\omega\,\Bigl\{\frac{e^{\lambda\omega} + e^{-\lambda\omega}}{e^{\lambda\omega} - e^{-\lambda\omega}}\,\text{sin}\,\lambda\omega - \text{cos}\,\lambda\omega\Bigr\} = 0
\end{align*}
\ldots\ldots\ldots\ldots\ldots\ldots\ldots\ldots\ldots\ldots\ldots\ldots libres § 8
\begin{align*}
&\Bigl\{\frac{2}{e^{\lambda\omega} + e^{-\lambda\omega}} + \text{cos}\,\lambda\omega\Bigr\}\Bigl\{\frac{e^{(1-\lambda)\omega} + e^{(\lambda-1)\omega}}{e^{(1-\lambda)\omega} - e^{(\lambda-1)\omega}}\,\text{sin}\,(1-\lambda)\omega - \text{cos}\,(1-\lambda)\omega\Bigr\} \\
&+ \Bigl\{\frac{2}{e^{(1-\lambda)\omega} + e^{(\lambda-1)\omega}} + \text{cos}\,(1-\lambda)\omega\Bigr\}\Bigl\{\frac{e^{\lambda\omega} + e^{-\lambda\omega}}{e^{\lambda\omega} - e^{-\lambda\omega}}\,\text{sin}\,\lambda\omega - \text{cos}\,\lambda\omega\Bigr\} = 0
\end{align*}
\ldots\ldots\ldots\ldots\ldots\ldots\ldots\ldots\ldots\ldots\ldots\ldots fixées § 12
\begin{align*}
&\Bigl\{\frac{2}{e^{\lambda\omega} + e^{-\lambda\omega}} - \text{cos}\,\lambda\omega\Bigr\}\Bigl\{\frac{e^{(1-\lambda)\omega} + e^{(\lambda-1)\omega}}{e^{(1-\lambda)\omega} - e^{(\lambda-1)\omega}}\,\text{sin}\,(1-\lambda)\omega - \text{cos}\,(1-\lambda)\omega\Bigr\} \\
&+ \Bigl\{\frac{2}{e^{(1-\lambda)\omega} + e^{(\lambda-1)\omega}} - \text{cos}\,(1-\lambda)\omega\Bigr\}\Bigl\{\frac{e^{\lambda\omega} + e^{-\lambda\omega}}{e^{\lambda\omega} - e^{-\lambda\omega}}\,\text{sin}\,\lambda\omega - \text{cos}\,\lambda\omega\Bigr\} = 0
\end{align*}
\ldots\ldots\ldots\ldots d'une extremité fixée et l'autre appuyée
\begin{align*}
&\Bigl\{\text{tang}\,\lambda\omega - \frac{e^{\lambda\omega} - e^{-\lambda\omega}}{e^{\lambda\omega} + e^{-\lambda\omega}}\Bigr\}\Bigl\{\text{tang}\,(1-\lambda)\omega - \frac{e^{(1-\lambda)\omega} - e^{(\lambda-1)\omega}}{e^{(1-\lambda)\omega} + e^{(\lambda-1)\omega}}\Bigr\} \\
&= 2\,\text{tang}\,\lambda\omega\,\frac{e^{\lambda\omega} - e^{-\lambda\omega}}{e^{\lambda\omega} + e^{-\lambda\omega}}\Bigl\{1 - \frac{2}{(e^{(1-\lambda)\omega} + e^{(\lambda-1)\omega})\,\text{cos}\,(1-\lambda)\omega}\Bigr\}
\end{align*}

Les équations qui déterminent les valeurs de l'angle $\omega$ lorsqu'on a simplement
égard à l'état des extremités, sont extraites du mémoire d'Euler; elles sont:

Pour le cas des deux extremités libres §. 5 \ldots\ldots $\text{cos}\,\omega = \frac{2e^{\omega}}{e^{2\omega} + 1}$

\ldots\ldots d'une extremité libre et l'autre appuyée §. 23 \ldots{} $\text{tang}\,\omega = \frac{e^{2\omega} - 1}{e^{2\omega} + 1}$

\ldots\ldots d'une extremité libre et l'autre fixée § 30 \ldots{} $\text{cos}\,\omega = -\frac{2e^{\omega}}{e^{2\omega} + 1}$

\ldots\ldots\ldots\ldots\ldots\ldots appuyées § 35 \ldots\ldots\ldots $\text{sin}\,\omega = 0$

\ldots\ldots d'une extremité appuyée et l'autre fixée §. 39 \ldots{} $\text{tang}\,\omega = \frac{e^{2\omega} - 1}{e^{2\omega} + 1}$

\ldots\ldots\ldots\ldots\ldots\ldots fixées § 43 \ldots\ldots{} $\text{cos}\,\omega = \frac{2e^{\omega}}{e^{2\omega} + 1}$
\note{En tête, au milieu, « (20) » barré ; en haut à droite « 175 », à
l'encre, le 5 tracé en boucle comme un S. Le paragraphe 16 porte son
numéro dans la ligne, le 1 réduit à un point épais. Le texte suit la vue
356. Les renvois « § 5 », « § 8 », « § 12 » qui
suivent les équations finales désignent, semble-t-il, ses propres
paragraphes : l'équation finale des extrémités fixées est celle du § 12
(vue 348), et celle des extrémités libres est établie juste avant le § 9
(vue 342). Les renvois de la liste finale, « extraites du mémoire
d'Euler », sont ceux des paragraphes d'Euler. Dans ces nombres, le 3 est
tracé comme un 9 à queue courte et le 9 a une longue queue : on lit 23,
30, 35, 39, 43, suivant ce que le passage sur la copie d'Euler a établi
pour cette forme du 3. Le premier renvoi se lit « §. 5 », comme « cas I
§ 5 » à la vue 344 ; la liste correspondante du brouillon, à la vue 276,
a été lue « N\textsuperscript{o} 9 ». Dans la
première équation finale, l'accolade ouverte après
$\text{sin}\,\lambda\omega$ n'est pas refermée avant « $+$ ». Dans la
dernière équation finale, le signe entre les deux exponentielles du
dernier dénominateur est surchargé ; le « 1 » de
$\text{cos}\,(1-\lambda)\omega$ y est empâté. La vue 359 reproduit ce
même feuillet, photographié une seconde fois : même texte, mêmes taches,
même « 175 » ; elle n'est pas transcrite.}

\end{document}
